\subsection{性高潮后疾病综合征（POIS）与射精后头痛}
\label{subsec:pois}

\noindent\textit{【编者注】}多数人在性高潮后感到放松与愉悦，但少数男性每次射精（含自慰、梦遗）后都会出现一组刻板的躯体与情绪不适。这类"性高潮后"问题长期被忽视，近年来才被逐渐识别与命名。

\vspace{0.5em}
\noindent\textbf{性高潮后疾病综合征（Post-Orgasmic Illness Syndrome, POIS）}：

\begin{itemize}[leftmargin=2cm]
	\item \textbf{表现}：射精后数分钟至数小时内出现一组固定症状，可持续 2--7 天，包括\textbf{全身疲惫、肌肉酸痛、发热感或"流感样"症状、注意力难以集中、情绪低落或易怒、鼻塞流涕、眼痒}等
	\item \textbf{特点}：症状与射精高度相关、每次发作模式相似、随休息自行缓解；可发生于任何年龄，但多见于青壮年
	\item \textbf{可能机制}：目前认为可能与\textbf{对自身精液的过敏/自身免疫反应}、神经内分泌（阿片--催乳素系统）波动、细胞因子介导的炎症反应等有关，尚无定论
	\item \textbf{诊断}：为排他性诊断，需先排除前列腺炎、慢性盆腔疼痛、抑郁焦虑、甲状腺疾病、感染等；建议就诊男科/泌尿外科或变态反应科
	\item \textbf{处理}：避免过度频繁射精、充分休息；有报道抗组胺药、非甾体抗炎药、脱敏治疗、选择性 5-羟色胺再摄取抑制剂等有一定效果，均需在医生指导下个体化尝试
\end{itemize}

\vspace{0.5em}
\noindent\textbf{射精后头痛（性活动相关性头痛）}：

\begin{itemize}[leftmargin=2cm]
	\item 在性兴奋或性高潮前后出现头痛，分为性高潮前的钝胀型与性高潮瞬间的爆发型
	\item 多数为良性，但\textbf{首次发作或突发剧烈"炸裂样"头痛}必须立即就医，排除蛛网膜下腔出血、脑动脉瘤等危及生命的病因
	\item 反复发作且已排除器质性病变者，可由神经内科评估预防性用药（非甾体抗炎药、$\beta$ 受体阻滞剂等）
\end{itemize}

\vspace{0.5em}
\noindent\textbf{何时就医}：

\begin{itemize}[leftmargin=2cm]
	\item 每次射精后都出现固定不适、持续超过 24 小时，或明显影响工作与生活
	\item 突发剧烈头痛，伴呕吐、颈项强直、意识改变或肢体无力
	\item 伴发热、血精、排尿疼痛、睾丸疼痛等提示感染或其他疾病的症状
\end{itemize}

\begin{tcolorbox}[colback=pink!5!white,colframe=red!50!black,title=贴心小叮咛]
射精后不适长期被视为"房事过度""肾虚"而饱受误解。现代医学提示，其中一部分是真实的\textbf{性高潮后疾病综合征}——它不是"意志薄弱"，而是一种可被识别、可干预的医学状况。若你每次射精后都规律性地出现不适，请带着"时间—症状—诱因"的记录就诊，而不是独自硬扛。
\end{tcolorbox}
